% === Lab Report Header and InfoBox ===
% Firefly returns in high energy: a sincere love letter to Firefly from Shuyun Yang!
% I used to wonder why “the moon is beautiful tonight” is said to be a subtle confession.
% Until I saw the golden scales shimmering under the morning sun on the river,
% or the reluctant orange clouds in the night sky,
% or the smoky scent of grilled skewers mingled with cheerful conversation at a late-night food stand —
% and instinctively wanted to share it with you.
% I want to pass on the things I find beautiful.
% Today is still sunny, my dear Firefly.

% ============================================================================
% Package Loading
\usepackage{amsmath}
\usepackage{amssymb}
\usepackage{graphicx}
\usepackage{xcolor}
\usepackage{enumerate}
\usepackage{multirow}
\usepackage{float}
\usepackage{cleveref}
\usepackage{adjustbox}
\usepackage{physics}
\usepackage{subfig}
\usepackage{tikz}
\usepackage{hyperref}
\usepackage{tabularx}
\usepackage{tabularray}
\usepackage{lipsum}
\usepackage{siunitx} % 规范显示单位与物理量 v1.1


% ============================================================================
% Experiment Info Box
\newcommand{\infoTable}[7]{
	\begin{flushleft}
		\Huge \textcolor{StructureColor}{\textbf{#1}}
	\end{flushleft}
	
	\vspace{-0.3cm}
	
	\begin{table}[h!]
		\textnormal{\textcolor{StructureColor}{\rule{0.75\textwidth}{1.53pt} }}\\
		\begin{tabularx}{0.7\textwidth}{p{0.175\textwidth}p{0.525\textwidth}}
			\textcolor{StructureColor}{\textbf{Date and Time:}} &  \textbf{#2}\\
			\textcolor{StructureColor}{\textbf{Location:}} &  \textbf{#3}\\
			\textcolor{StructureColor}{\textbf{Environment:}} &  \textbf{#4} \\
			\textcolor{StructureColor}{\textbf{Student 1:}} &  \textbf{#5} \\
			#6
			\textcolor{StructureColor}{\textbf{Instructor(s):}} &  \textbf{#7}\\
		\end{tabularx}\\
		\textnormal{\textcolor{StructureColor}{\rule{0.75\textwidth}{1.5pt} }}
	\end{table}
	
	% Notes
	\noindent\textcolor{StructureColor}{\textbf{Notes}}
	
	The lab report consists of three main parts: \textbf{Preparation}, \textbf{Experiment}, and \textbf{Analysis \& Discussion}, along with a cover page and appendix.
	
	The preparation report must be completed before the lab session, including a thorough review of the manual, understanding of the principles, familiarity with instruments, and completion of assigned questions. A pre-designed data recording sheet is recommended.
	
	Experimental records must be written in pen, signed, and kept complete — \textbf{pencil records are invalid}. Any mistakes should be corrected according to standard lab procedures, and the original content must remain visible.
	
	Data analysis must involve real calculations (unless the lab is instrument-learning based), error estimates, and formatted tables and graphs with proper captions and references.
	
	Reports must be submitted within one week of completing the lab (or two weeks in special cases).
	
	\noindent\textcolor{StructureColor}{\rule{\textwidth}{1.5pt} }
	
	% License Notice
	\noindent\textcolor{StructureColor}{\textbf{Special Notice}}
	
	\textbf{This template is distributed under the MIT License.} By using it, you agree to comply with the associated terms.
	
	This template was developed by GitHub user \textbf{pifuyuini}, based on a 2019 version from another Sun Yat-sen University student. The visual layout was inspired by \href{https://www.yykspace.com/cn/index.html}{a prominent lab report template}. Thanks to both contributors (even if they’re unaware).
	
	As we transition to advanced physics experiments, we move beyond rigid templates toward writing in the style of academic papers — this template reflects that goal. If this diverges from your instructor's expectations, we kindly ask for your understanding.
}

% Emphasis
\newcommand{\YsyEmph}[1]{
	\textbf{\textcolor{EmphColor}{#1}}
}

% ============================================================================
% 补充设定
% 规范显示单位与物理量 v1.1
\sisetup{
	per-mode=symbol,      % 使用 '/' 而非 '·'
	separate-uncertainty = true,  % 不确定度写作 ± 形式
	detect-all = true,    % 自动检测字体样式（math, text等）
	inter-unit-product = \cdot,
}

% ============================================================================
% 封装

% 物理
\newcommand{\tensor}[3]{#1^{\mathrm{#2}}_{\phantom{#2}\mathrm{#3}}} % 张量

% 排版规范 v1.1
% 参考：https://www.bilibili.com/video/BV1WrdkY3ECV
\newcommand{\Diff}{\mathrm{d}} % 微分 直立体
\newcommand{\Der}[2]{\dfrac{\mathrm{d}#1}{\mathrm{d}#2}} % 导数
\newcommand{\EulerE}{\mathrm{e}} % 自然常数 正体
\newcommand{\ImgI}{\mathrm{i}} % 虚数单位 正体
\newcommand{\DV}[2]{#1_{\mathrm{#2}}} % 对偶矢量 角标是正体
\newcommand{\LA}[1]{\boldsymbol{#1}} % 对偶矢量 角标是正体
